\documentclass[11pt,a4paper]{article}
\usepackage[a4paper,margin=2.35cm]{geometry}
% Overleaf-compatible font setup. The document can be compiled with
% pdfLaTeX, XeLaTeX, or LuaLaTeX; no proprietary system fonts are required.
\usepackage{iftex}
\ifPDFTeX
  \usepackage[T1]{fontenc}
  \usepackage[utf8]{inputenc}
  \usepackage{newtxtext}
  \usepackage[scaled=0.92]{helvet}
\else
  \usepackage{fontspec}
  \setmainfont{TeX Gyre Termes}
  \setsansfont{TeX Gyre Heros}[Scale=0.92]
\fi
\usepackage[nopatch=footnote]{microtype}
\usepackage{amsmath,amssymb,mathtools}
\ifPDFTeX
  \usepackage{newtxmath}
\fi
\usepackage{booktabs,longtable,array,tabularx}
\newcolumntype{Y}{>{\RaggedRight\arraybackslash}X}
\usepackage{enumitem}
\usepackage{xcolor}
\usepackage[round,authoryear]{natbib}
\usepackage{hyperref}
\usepackage{fancyhdr}
\usepackage{titlesec}
\usepackage{caption}
\usepackage{float}
\usepackage{hanging}
\usepackage{ragged2e}
\usepackage{setspace}

\definecolor{navy}{HTML}{17324D}
\definecolor{midgray}{HTML}{5F6870}
\titleformat{\section}{\Large\bfseries\color{navy}}{\thesection.}{0.6em}{}
\titleformat{\subsection}{\large\bfseries\color{navy}}{\thesubsection}{0.6em}{}
\titleformat{\subsubsection}{\normalsize\bfseries\color{navy}}{\thesubsubsection}{0.6em}{}
\titlespacing*{\section}{0pt}{1.3em}{0.55em}
\titlespacing*{\subsection}{0pt}{1.0em}{0.4em}
\titlespacing*{\subsubsection}{0pt}{0.8em}{0.3em}
\setlist{leftmargin=1.6em,itemsep=0.3em,topsep=0.4em}
\renewcommand{\arraystretch}{1.2}
\setlength{\parindent}{0pt}
\setlength{\parskip}{0.55em}
\onehalfspacing

\pagestyle{fancy}
\fancyhf{}
\fancyhead[R]{\small\color{midgray}Research Proposal}
\fancyfoot[C]{\thepage}
\setlength{\headheight}{14pt}

\setlength{\emergencystretch}{2em}
\setlength{\bibhang}{1.5em}
\setlength{\bibsep}{0.6em}
\renewcommand{\bibfont}{\small\RaggedRight}
\urlstyle{same}
\Urlmuskip=0mu plus 1mu
\captionsetup[table]{font=small,labelfont=bf,skip=6pt}
\setlength{\LTpre}{0.6em}
\setlength{\LTpost}{0.6em}
\newcommand{\term}[1]{\textit{#1}}
\hypersetup{colorlinks=true,linkcolor=navy,urlcolor=navy,citecolor=navy,pdfauthor={Xiangyu Xu},pdftitle={How Does AI Change Information Acquisition and Opinion Formation in Human Groups?}}
\fancyhead[L]{\small\color{midgray}AI and Group Opinion Formation}

\begin{document}

\begin{titlepage}
\centering
\vspace*{2.5cm}
{\sffamily\bfseries\color{navy}\LARGE Research Proposal\par}
\vspace{1.2cm}
{\sffamily\bfseries\Huge How Does AI Change Information Acquisition and Opinion Formation in Human Groups?\par}
\vspace{1.5cm}
{\sffamily\Large Xiangyu Xu\par}
\vspace{0.8cm}

\vfill
\end{titlepage}

\section*{Abstract}

People who use the same large language model may receive similar information. Its effect on group opinions depends on what members read, share, and accept. This proposal compares traditional search with LLM information search. It then examines how common exposure and member sharing affect opinion change. A hybrid agent-based model combines explicit behavioral rules with LLM agents. Human data will calibrate the model and test its predictions on held-out observations. The analysis allows convergence, persistent disagreement, and clustering. If a stable group-level pattern emerges, a further comparison will test independent routing across model families. A shared human-curated brief will assess whether common information alone accounts for the result. The study aims to connect the information supplied by LLMs to the development of group opinions.

\textbf{Keywords:} large language models; information search; information diffusion; opinion dynamics; agent-based modeling

\section{Introduction}

LLMs generate explanations and arguments that users can read, restate, and share. Separate conversations may expose different users to similar claims. Shared information can then spread through a group. Whether this process brings opinions closer together depends on how members select and use that information.

\citet{jiang2025} found repetition within language models and homogeneity across models in open-ended generation. Related studies found that AI assistance can improve individual performance while increasing correlations among choices or reducing the diversity of creative output~\citep{fugener2021,doshi2024}. These findings concern model output and assisted performance. The effect on human opinion change remains an empirical question.

Information exposure and opinion change can diverge. \citet{nyhan2023} reduced exposure to like-minded sources on Facebook and found no measurable change in eight preregistered political-attitude outcomes. Models of algorithmic personalization also show that filtering and network structure can produce different opinion distributions~\citep{perra2019}. The present study examines the steps between receiving information and changing an opinion.

The main question is: \emph{How do large language models change group opinion distributions?} Three subquestions guide the analysis:

\begin{enumerate}
\item How does LLM information search change exposure and diffusion compared with traditional search?
\item How do exposure and diffusion affect individual opinion updating and the group distribution?
\item If a stable pattern is observed, can independent routing across model families attenuate it?
\end{enumerate}

A hybrid agent-based model (ABM) will trace messages, exposure, sharing, and opinion change. Human data will estimate behavioral parameters and provide independent validation. The model concerns groups connected by information-sharing networks. Members have no formal team task or collective decision rule.

\section{Theory and Research Questions}

\subsection{Shared information and selective exposure}

Shared use of an LLM may make members' information more similar. This expectation follows from evidence on model-output homogeneity~\citep{jiang2025}. Algorithmic monoculture describes the use of the same algorithm by multiple decision-makers~\citep{kleinberg2021}. Its relevance here concerns correlations in the information that members receive. A model shared by the group may repeatedly supply similar arguments, sources, and errors.

Personalization can produce a different pattern. A model may generate arguments consistent with each user's initial position. Members may then receive different accounts of the same issue. \citet{chuang2024} found that confirmation-bias prompts could produce fragmentation in groups of LLM agents. This simulation finding motivates a comparison between neutral prompts and prompts that include the user's initial position. Neutral prompts remain the main setting; stance-conditioned prompts provide a supplementary comparison.

Group information research also distinguishes possession from sharing. Shared information is more likely to enter group discussion than information held by one member~\citep{stasser1985}. Members may reveal, conceal, or misrepresent what they know~\citep{steinel2010}. The information initially supplied to a member therefore gives only part of the account. Sharing determines what reaches other members and how often they encounter it.

SQ1 compares traditional search-engine use with LLM information search. It measures message content, actual exposure, and subsequent diffusion. This comparison will establish whether differences in generated information persist in what members actually encounter.

\subsection{Opinion updating}

Opinion-dynamics models explain how individual influence produces group distributions. DeGroot's model represents updating through weighted social influence~\citep{degroot1974}. The Friedkin--Johnsen model retains attachment to initial opinions~\citep{friedkin1990}. Bounded-confidence models restrict influence to sufficiently close opinions~\citep{hegselmann2002}. These models allow shared information to produce different outcomes under different updating rules.

Common exposure may bring opinions closer together when members accept similar arguments. Selective sharing and limited acceptance may sustain separate clusters. The same information intervention may also leave opinions largely unchanged. SQ2 tests these possibilities by varying common exposure and diffusion while holding the message library and initial groups fixed. The primary distributional outcome and directional predictions remain to be selected before confirmatory analysis.

An opinion is a member's position on a public-policy issue. It will be measured on a common opposition--support scale. The main outcomes describe the distribution of these positions. Convergence and polarization carry no automatic implication about judgment quality. Factual accuracy will be assessed separately for claims with verifiable truth values.

\subsection{Routing across model families}

Algorithmic monoculture motivates a further comparison between a shared model and independently assigned models~\citep{kleinberg2021}. Different model families may still produce similar content. SQ3 therefore tests whether routing changes actual information exposure and then attenuates the same group outcome observed in SQ2.

Each member will draw a model independently from a specified pool at the start of an issue. The member will retain that model for the issue. Routing probabilities will be fixed in advance. Assignment will be independent of initial opinion, prompt content, difficulty, cost, and predicted response quality. Only the information-source models will vary. Models used to simulate people will remain fixed.

\section{Research Design}

\subsection{Setting and information conditions}

The task will ask members to consider public-policy issues on a common support scale. Approximately 24 issues is a candidate scope. The target population, issue set, recruitment or existing-data source, search platform, and empirical network remain to be specified. The meaning of one simulation round will follow the selected human task.

The main source comparison uses traditional search engines and a shared LLM. Human search records must capture queries, returned sources, selected materials, and actual exposure. LLM records will retain prompts, responses, and follow-up exchanges. The routing condition described above will be used for SQ3.

A shared human-curated brief will provide an alternative-source check. If it reproduces the LLM effect under comparable content quality and diffusion, the result supports a common-information explanation. Claims about AI specificity require evidence beyond a difference between model labels.

Conditions will approximately match information quantity, length, factual accuracy, mean stance, argument balance, and readability where these are controlled. Such matching defines the comparison being estimated. It leaves some naturally occurring differences between technologies outside that comparison.

\subsection{Hybrid model}

Rule-based agents and LLM agents will operate within the same population and network. They will share an opinion scale, state variables, and information records. Rule-based agents follow explicit equations. LLM agents generate text and opinion responses from individual profiles, current opinions, memory, and received information. HiSim provides a precedent for combining LLM-driven users with deductive ABM agents~\citep{mou2024}.

The model separates LLM information sources from LLMs used to simulate people. It will record the model, prompt, input, and output for each role. Outcome coding will be kept separate from information generation and human simulation.

Three interfaces define agent behavior:

\begin{equation*}
m_{j,t}=f_{\text{message}}(x_{j,t},z_j,H_{j,t},E_t),
\end{equation*}

\begin{equation*}
J_{i,t}=f_{\text{selection}}(i,G_t,\{m_{j,t}\},x_{i,t}),
\end{equation*}

\begin{equation*}
x_{i,t+1}=f_{\text{update}}
\left(x_{i,t},\{m_{j,t}:j\in J_{i,t}\},z_i,H_{i,t}\right).
\end{equation*}

Here, $x_{i,t}$ is member $i$'s opinion in round $t$. Individual characteristics are denoted by $z_i$, memory by $H_{i,t}$, the external environment by $E_t$, and the network by $G_t$. Memory retains the initial opinion. The set $J_{i,t}$ contains identifiers of messages actually encountered by member $i$.

The message function produces transmissible content. Rule-based agents use templates or numerical mappings; LLM agents generate natural-language messages. The selection function determines actual exposure. Both agent types use the same explicit selection rule in the main model. The update function returns the next opinion on the common scale.

Onward sharing is controlled separately. Human-calibrated sharing probabilities govern diffusion in the main condition. Additional sharing suggestions from LLM agents will be recorded without overriding this rule. Each message and sharing event will be executed once. Its production and exposure rounds will be logged. The exact within-round schedule and any supplementary LLM-sharing condition remain to be fixed.

\subsection{Message content and exposure}

Each message contains coded stance, claims, and sources:

\begin{equation*}
m_{j,t}
=
\left(
s_{j,t},
\mathbf q_{j,t},
\mathbf r_{j,t}
\right),
\end{equation*}

The stance $s_{j,t}$ lies in $[-1,1]$. The vector $\mathbf q_{j,t}$ records the presence of each coded claim or argument. The vector $\mathbf r_{j,t}$ records the information sources used. Each message also has a sender and a source identifier. These records allow messages, claims, and sources to be tracked separately.

Coding will distinguish wording similarity, semantic similarity, claim overlap, source overlap, and stance. Automated coding will be checked against human annotations. Error correlation will be calculated only for factual claims with independently assessed truth values. Policy-support positions remain opinion outcomes.

Let $Q_{i,t}$ contain the claims in candidate messages available to member $i$ before selection. Candidate-message overlap is:

\begin{equation*}
C^{M}_{g,t}
=
\frac{2}{n(n-1)}
\sum_{i<j}
\frac{|Q_{i,t}\cap Q_{j,t}|}
{|Q_{i,t}\cup Q_{j,t}|},
\end{equation*}

Actual exposure is measured using $K^{E}_{i,t}=\{k:\exists j\in J_{i,t},\ q_{j,t,k}=1\}$, the claims contained in encountered messages:

\begin{equation*}
C^{E}_{g,t}
=
\frac{2}{n(n-1)}
\sum_{i<j}
\frac{|K^{E}_{i,t}\cap K^{E}_{j,t}|}
{|K^{E}_{i,t}\cup K^{E}_{j,t}|}.
\end{equation*}

Both measures average pairwise Jaccard overlap within a group of $n$ members. Empty unions are undefined and will be reported as unavailable. Source overlap and exact-message overlap will use their own identifiers. Repeated exposure to claim $k$ is measured as the number of rounds in which it appears in member $i$'s exposure set:

\begin{equation*}
R_{i,k,t}
=
\sum_{\tau=1}^{t}
\mathbb I(k\in K^{E}_{i,\tau}).
\end{equation*}

\subsection{Selection and diffusion}

The candidate selection rule is:

\begin{equation*}
P(j\in J_{i,t})
=
\operatorname{logit}^{-1}
\left[
\beta_0
+
\beta_d|s_{j,t}-x_{i,t}|
+
\beta_rR_{i,j,t}
+
\beta_nA_{ij}
\right].
\end{equation*}

The intercept $\beta_0$ is baseline log-odds of exposure. Message--opinion distance enters through $\beta_d$. The term $R_{i,j,t}$ counts prior rounds of exposure to the source associated with message $j$. Its coefficient is $\beta_r$. The indicator $A_{ij}$ records a network connection between member $i$ and that source. If every candidate source is connected, $A_{ij}$ is constant and its coefficient is omitted from estimation.

A separate candidate rule governs sharing:

\begin{equation*}
P(Y_{i,j,t}^{\text{share}}=1)
=
\operatorname{logit}^{-1}
\left[
\gamma_0
+
\gamma_d|s_{j,t}-x_{i,t}|
+
\gamma_rR_{i,j,t}
+
\gamma_sS_j
\right],
\end{equation*}

Here, $\gamma_0$ is baseline log-odds of sharing. The remaining terms represent message--opinion distance, prior source exposure, and source category $S_j$. The coefficients will be estimated from human selection and sharing data or assigned explicit sensitivity ranges. They remain fixed across source conditions in the main simulation. Differences in estimated human coefficients require condition-specific human observations.

Repeated sharing generates diffusion paths. The analysis will report sharing frequency, recipients, diffusion depth and breadth, and the proportion of diffusion across opinion positions. Each event will remain linked to the message and claims transmitted.

\subsection{Opinion updating}

A candidate rule for numerical agents follows the Friedkin--Johnsen structure:

\begin{equation*}
x_{i,t+1}
=
(1-\lambda_i)x_{i,0}
+
\lambda_i
\sum_{j\in J_{i,t}}
w_{ij,t}s_{j,t},
\qquad
\sum_jw_{ij,t}=1.
\end{equation*}

The term $x_{i,0}$ is the initial opinion. The parameter $\lambda_i$ gives the weight of social influence, and $w_{ij,t}$ gives the relative weight of message $j$. The equation adapts the original model's social inputs to coded message stances~\citep{friedkin1990}. A bounded-confidence alternative restricts which messages receive weight:

\begin{equation*}
w_{ij,t}=0
\quad\text{if}\quad
|s_{j,t}-x_{i,t}|>\varepsilon_i,
\end{equation*}

The threshold $\varepsilon_i$ corresponds to acceptable opinion distance~\citep{hegselmann2002}. Weights are normalized over acceptable messages. The state transition for an empty acceptable set remains an implementation choice and will be fixed before validation.

LLM agents receive the same exposure bundle and return a next-period opinion. Their numerical response and stated rationale will be retained for consistency checks. Their transitions will be compared with human opinion changes under equivalent information conditions. The final rule for numerical agents will be selected using held-out model-selection data.

\subsection{Human calibration and validation}

Human observations will include initial opinions, controlled information exposure, subsequent opinions, and sharing choices. These data will estimate opinion inertia, acceptance by opinion distance, repeated-exposure effects, and parameter variation across people and issues. Agent profiles will use characteristics supported by these observations.

Development data will support parameter estimation and prompt construction. Separate data will support model selection, followed by an independent final validation sample. The sample sources and split sizes remain to be specified. Sample size will follow simulation-based power analysis using the group design and a minimum effect of interest.

Validation will compare message content, selection, sharing, and opinion transitions with human observations. A rule-only ABM, the hybrid ABM, and simple baselines will receive the same held-out cases. The comparison will assess whether LLM agents improve prediction or provide validated records of information processes. Fluent explanations alone provide insufficient evidence of human behavior.

Group trajectories will also be compared when suitable human group data are available. Without those data, long-run group outcomes remain empirically constrained simulations. Validation will use several patterns, including opinion means, dispersion, and diffusion paths~\citep{windrum2007,grimm2005}.

Paired runs will use the same initial group and network, with common random numbers. Sensitivity checks will vary the agent mix, network, update rule, model version, and prompt. Program checks will cover empty exposure, identical and opposing messages, and disconnected networks. Prompts, model versions, failed calls, and random seeds will be recorded. The implementation will be documented using the ODD protocol~\citep{grimm2020}.

\section{Comparisons and Analysis}

\subsection{SQ1: Information sources}

The source comparison will examine message stance, claim and source occurrence, candidate-message overlap $C^M$, and actual-exposure overlap $C^E$. It will then compare sharing and diffusion paths. Search behavior must be represented by corresponding human records. Fixed behavioral coefficients remain model inputs throughout this comparison.

The analysis will account for issue effects and repeated responses, members, and networks. It will report content, exposure, and sharing as distinct outcomes. The primary statistical unit is the group--issue combination. Random seeds quantify simulation uncertainty and will not be counted as independent human observations.

\subsection{SQ2: Exposure and diffusion}

SQ2 uses a fixed message library and identical initial groups and networks. The selection rule varies exposure commonality $c\in\{c_{\text{low}},c_{\text{high}}\}$. High commonality assigns more shared claims; low commonality assigns more distinct claims. Actual $C^E$ will check the manipulation.

The sharing rule varies diffusion intensity $\pi\in\{0,\hat\pi_{\text{human}}\}$. At zero, members cannot pass information onward. At $\hat\pi_{\text{human}}$, human-calibrated probabilities govern sharing. The four conditions are shown in Table~\ref{tab:exposure_diffusion}.

{\small
\begin{longtable}{@{}>{\RaggedRight\arraybackslash}p{0.16\textwidth}>{\RaggedRight\arraybackslash}p{0.36\textwidth}>{\RaggedRight\arraybackslash}p{0.36\textwidth}@{}}
\caption{Exposure and diffusion conditions}\label{tab:exposure_diffusion}\\
\toprule
\textbf{} & \textbf{$\pi=0$} & \textbf{$\pi=\hat\pi_{\text{human}}$} \\
\midrule
\endfirsthead
\toprule
\textbf{} & \textbf{$\pi=0$} & \textbf{$\pi=\hat\pi_{\text{human}}$} \\
\midrule
\endhead
$c=c_{\text{high}}$ & High common exposure, no diffusion & High common exposure, diffusion enabled \\
$c=c_{\text{low}}$ & Low common exposure, no diffusion & Low common exposure, diffusion enabled \\
\bottomrule
\end{longtable}
}

Comparisons will estimate the separate effects of common exposure and diffusion, together with their interaction. Message quantity, quality, mean stance, and readability will be approximately matched. For numerical agents, the analysis will examine variation across social-influence weights $\lambda_i$ and acceptance thresholds $\varepsilon_i$.

For LLM agents, opinion change is $\Delta x^{\text{LLM}}_{i,t}=x^{\text{LLM}}_{i,t+1}-x^{\text{LLM}}_{i,t}$. The candidate response model is:

\begin{equation*}
\Delta x^{\text{LLM}}_{i,t}
=
\theta_0
+
\theta_1(\bar s_{i,t}-x_{i,t})
+
\theta_2c
+
\theta_3\pi
+
\theta_4c(\bar s_{i,t}-x_{i,t})
+
u_{i,t}.
\end{equation*}

Here, $\bar s_{i,t}$ is the mean stance in the exposure bundle. The coefficient $\theta_1$ summarizes the relation between message--opinion distance and state change. It describes an observed response of the simulated agent. Human, LLM, and numerical agents will be compared on the direction and magnitude of opinion change. Their parameters have different definitions, so direct parameter-equality tests would be inappropriate.

Change in within-group opinion variance is a candidate primary outcome. Mean position, pairwise distance, clustering, and bimodality are supplementary candidates. The primary measure, comparison direction, and minimum effect of interest will be fixed before confirmatory runs. Uncertainty in human-calibrated parameters and generalization across issues will be included in the analysis.

\subsection{SQ3: Independent model routing}

SQ3 will proceed if SQ2 establishes a stable group-level pattern. It will compare a shared information-source model with independent assignment across model families or providers. Models used to simulate people, initial groups, and behavioral settings will stay fixed.

The test requires a change in a specified information measure and attenuation of the same group outcome used in SQ2. The common-information check using a human-curated brief will help interpret the result. Any proposed benefit of routing will be limited to the outcomes and conditions tested.

\section{Expected Contribution}

The study examines how LLM information enters group opinion formation through exposure and sharing. It connects evidence on model-output homogeneity to an observable sequence of messages, encounters, and opinion changes. This sequence can distinguish common content from common exposure and similar updating behavior.

The hybrid model retains explicit opinion rules and records the language transmitted between members. Its value will depend on validation against human behavior and comparison with a rule-only model. The routing comparison adds a test of whether changing information sources can alter an established group pattern. The resulting claims will concern the specific information processes and group outcomes supported by the data.

\begingroup
\RaggedRight
\begin{thebibliography}{99}
\RaggedRight
\interlinepenalty=10000

\bibitem[Chuang et al.(2024)]{chuang2024}
Chuang, Y.-S., Goyal, A., Harlalka, N., et al. (2024). Simulating opinion dynamics with networks of LLM-based agents. In \textit{Findings of the Association for Computational Linguistics: NAACL 2024} (pp. 3326--3346). \url{https://aclanthology.org/2024.findings-naacl.211/}

\bibitem[DeGroot(1974)]{degroot1974}
DeGroot, M. H. (1974). Reaching a consensus. \textit{Journal of the American Statistical Association, 69}(345), 118--121. \url{https://doi.org/10.1080/01621459.1974.10480137}

\bibitem[Doshi and Hauser(2024)]{doshi2024}
Doshi, A. R., \& Hauser, O. P. (2024). Generative AI enhances individual creativity but reduces the collective diversity of novel content. \textit{Science Advances, 10}(28). \url{https://doi.org/10.1126/sciadv.adn5290}

\bibitem[Friedkin and Johnsen(1990)]{friedkin1990}
Friedkin, N. E., \& Johnsen, E. C. (1990). Social influence and opinions. \textit{Journal of Mathematical Sociology, 15}(3--4), 193--206. \url{https://doi.org/10.1080/0022250X.1990.9990069}

\bibitem[F\"ugener et al.(2021)]{fugener2021}
F\"ugener, A., Grahl, J., Gupta, A., et al. (2021). Will humans-in-the-loop become Borgs? Merits and pitfalls of working with AI. \textit{MIS Quarterly, 45}(3), 1527--1556. \url{https://doi.org/10.25300/MISQ/2021/16553}

\bibitem[Grimm et al.(2005)]{grimm2005}
Grimm, V., Revilla, E., Berger, U., et al. (2005). Pattern-oriented modeling of agent-based complex systems: Lessons from ecology. \textit{Science, 310}(5750), 987--991. \url{https://doi.org/10.1126/science.1116681}

\bibitem[Grimm et al.(2020)]{grimm2020}
Grimm, V., Railsback, S. F., Vincenot, C. E., et al. (2020). The ODD protocol for describing agent-based and other simulation models: A second update to improve clarity, replication, and structural realism. \textit{Journal of Artificial Societies and Social Simulation, 23}(2). \url{https://doi.org/10.18564/jasss.4259}

\bibitem[Hegselmann and Krause(2002)]{hegselmann2002}
Hegselmann, R., \& Krause, U. (2002). Opinion dynamics and bounded confidence models, analysis, and simulation. \textit{Journal of Artificial Societies and Social Simulation, 5}(3). \url{https://www.jasss.org/5/3/2.html}

\bibitem[Jiang et al.(2025)]{jiang2025}
Jiang, L., Chai, Y., Li, M., et al. (2025). Artificial Hivemind: The open-ended homogeneity of language models. \textit{Advances in Neural Information Processing Systems}. \url{https://openreview.net/forum?id=saDOrrnNTz}

\bibitem[Kleinberg and Raghavan(2021)]{kleinberg2021}
Kleinberg, J., \& Raghavan, M. (2021). Algorithmic monoculture and social welfare. \textit{Proceedings of the National Academy of Sciences, 118}(22), e2018340118. \url{https://doi.org/10.1073/pnas.2018340118}

\bibitem[Mou et al.(2024)]{mou2024}
Mou, X., Wei, Z., \& Huang, X. (2024). Unveiling the truth and facilitating change: Towards agent-based large-scale social movement simulation. In \textit{Findings of the Association for Computational Linguistics: ACL 2024} (pp. 4789--4809). \url{https://doi.org/10.18653/v1/2024.findings-acl.285}

\bibitem[Nyhan et al.(2023)]{nyhan2023}
Nyhan, B., Settle, J., Thorson, E., et al. (2023). Like-minded sources on Facebook are prevalent but not polarizing. \textit{Nature, 620}, 137--144. \url{https://doi.org/10.1038/s41586-023-06297-w}

\bibitem[Perra and Rocha(2019)]{perra2019}
Perra, N., \& Rocha, L. E. C. (2019). Modelling opinion dynamics in the age of algorithmic personalisation. \textit{Scientific Reports, 9}, 7261. \url{https://doi.org/10.1038/s41598-019-43830-2}

\bibitem[Stasser and Titus(1985)]{stasser1985}
Stasser, G., \& Titus, W. (1985). Pooling of unshared information in group decision making: Biased information sampling during discussion. \textit{Journal of Personality and Social Psychology, 48}(6), 1467--1478. \url{https://doi.org/10.1037/0022-3514.48.6.1467}

\bibitem[Steinel et al.(2010)]{steinel2010}
Steinel, W., Utz, S., \& Koning, L. (2010). The good, the bad and the ugly thing to do when sharing information: Revealing, concealing and lying depend on social motivation, distribution and importance of information. \textit{Organizational Behavior and Human Decision Processes, 113}(2), 85--96. \url{https://doi.org/10.1016/j.obhdp.2010.07.001}

\bibitem[Windrum et al.(2007)]{windrum2007}
Windrum, P., Fagiolo, G., \& Moneta, A. (2007). Empirical validation of agent-based models: Alternatives and prospects. \textit{Journal of Artificial Societies and Social Simulation, 10}(2), 8. \url{https://www.jasss.org/10/2/8.html}

\end{thebibliography}
\endgroup

\end{document}
